\section{Why the nonlinear case does not converge, and what that means}
\label{sec:uv}

Section~\ref{sec:cutoff} reports that the tree-level vertex converges near
$\ell\sim10^4$ while the nonlinear one does not, and Section~\ref{sec:nonlinear}
reports that the nonlinear FK term reaches $41\%$ of Order-0 at the highest
cutoff tried. Phrased as ``not converged'', that invites the response that
the calculation was not carried far enough. It was not a matter of carrying
it further, and the distinction matters for what the draft should claim.

\subsection{Where the sensitivity comes from}

The FK term samples the driving-field three-point cumulant in the collapsed
configuration, with two of the three legs at the same point on the ray. A
cumulant with coincident legs is a variance-like object, and variances of
cosmological fields are dominated by the smallest scales retained. This is
the same statement as the nonlinear density variance depending on its
smoothing scale, and it is a property of the quantity rather than a defect
of the calculation.

After the Limber radial collapse the hard pair contributes as a
two-dimensional transverse integral $\int\!\mathrm{d}k\,k\,P(k)$, whose
integrand goes as $k^{1+n_{\rm eff}}$: it peaks where $n_{\rm eff}=-1$ and
converges once $n_{\rm eff}<-2$. Measured on the fiducial spectra at $z=1$:

\begin{center}
\begin{tabular}{rrrrr}
\toprule
$k$ [$h$/Mpc] & $n_{\rm eff}$ linear & $1+n_{\rm eff}$ & $n_{\rm eff}$ halofit & $1+n_{\rm eff}$ \\
\midrule
0.1 & $-1.67$ & $-0.67$ & $-1.60$ & $-0.60$ \\
1.0 & $-2.29$ & $-1.29$ & $-1.15$ & $-0.15$ \\
10  & $-2.61$ & $-1.61$ & $-2.08$ & $-1.08$ \\
100 & $-2.77$ & $-1.77$ & $-2.39$ & $-1.39$ \\
\bottomrule
\end{tabular}
\end{center}

The linear spectrum crosses $-2$ at $k=0.21\,h$/Mpc, so the tree-level
integral has a finite limit: $50$, $90$ and $99\%$ of it lie below $k=0.61$,
$8.6$ and $85\,h$/Mpc. The one-halo term flattens the nonlinear spectrum so
that it crosses $-2$ only at $k=8.4\,h$/Mpc, and at $k=1$ it sits at
$n_{\rm eff}=-1.15$, which is essentially the peak of the integrand: every
decade then contributes comparably, and $50$, $90$ and $99\%$ of the
integral lie below $k=9.6$, $80$ and $179\,h$/Mpc.

\subsection{It is not the Limber approximation}

Two things point the same way. The coincident-leg structure exists before
any Limber approximation is made, so the sensitivity is not introduced by
it. And Limber replaces a three-dimensional hard-mode integral by a
two-dimensional transverse one, relaxing the convergence condition from
$n_{\rm eff}<-3$ to $n_{\rm eff}<-2$; dropping Limber would therefore make
the divergence worse, not better. Limber's own failure mode is at low
multipole and wide angle, which is a different regime from this one.

\subsection{What is actually missing}

The formalism as written carries no smoothing scale. A search of
\code{sections/} finds no beam width, no coarse-graining, and $\elmax$ does
not appear in the text at all: it exists only in the code. So the quantity
being computed is a coincident-point moment of a mathematically point-like
beam, and that quantity has no finite value for a nonlinear input. The
cutoff is acting as an unstated regulator that sets the answer.

The tree-level case escapes this because it happens to be ultraviolet safe,
so no regulator is needed and the number is well defined. That is the
regime the draft is in, and nothing here undermines it. What does not
follow is that the same calculation can be upgraded to a nonlinear input by
substituting a better bispectrum.

\subsection{A finite beam would supply the missing scale}

The physical candidate is that a real beam has finite cross-section, so the
driving fields entering the Sachs system are averages over that section
rather than point values. The algebra of that proposal is worked out in
\code{SFT-lensing-paper-analyses/mathematica/beam\_width\_regulator.wl},
15 checks, all passing. What it establishes:

\begin{itemize}
\item For a normalised axisymmetric screen-plane window of width $s_b$,
averaging gives
$\bar\Phi = \Phi + (s_b^2/4)\,\nabla^2_\perp\Phi + O(s_b^4)$, and in Fourier
space it is multiplication by $\tilde W(k)=\exp(-k^2s_b^2/4)$, whose
small-$k$ expansion reproduces the same coefficient.
\item An axisymmetric window does not mix spin weights, so $\Phi_{00}$ stays
spin 0 and $\Psi_0$ stays spin 2 under averaging. The projection of
$e^{is\theta}$ against such a window vanishes for every $s\ne0$.
\item The windowed transverse integral converges for \emph{any} power-law
spectrum, because $\tilde W^2$ decays faster than any power. At the halofit
slope $n_{\rm eff}=-1.15$ the unwindowed integral diverges and the windowed
one is finite.
\item The window suppresses modes above $k\sim2/s_b$, so with the Limber map
it fixes $\elmax = 2\chi_h/s_b$. A beam of $1\,h^{-1}$Mpc gives
$\elmax\simeq590$ at the innermost shell and $11\,000$ at the source shell;
a beam of $3\,h^{-1}$Mpc gives $200$ and $3700$.
\end{itemize}

So the proposal does what is needed: it removes the divergence, and it
replaces a free cutoff by a scale that is a property of the observation.

\subsection{What the script does not settle}

Two limitations, both recorded in the script rather than left implicit.

It does not establish that beam averaging is the correct definition of this
observable, nor which scale to use. Candidates differ by orders of
magnitude: the angular size of the source, the resolution element of the
instrument, and the smoothing applied to the convergence map are all
defensible readings and give very different $\elmax$. That is a choice
about what is being measured, and it has to be made before a nonlinear
number can be quoted.

And averaging the drivers is not the same operation as solving for a finite
beam. The Sachs system is nonlinear in the optical scalars, so averaging
does not commute with the quadratic terms; the bundle variance survives,
and the two differ at second order in the optical scalars. Beam averaging
therefore defines the source, and a genuine finite-beam treatment would be
a further step.

\subsection{Consequence for the draft}

The presentation that follows from this is not ``the FK amplitude is
$X$, subject to a cutoff''. It is:

\begin{itemize}
\item the tree-level FK term is well defined and converges, to about
$2.5\%$ of Order-0 at arcminute separations, five times the value the draft
quotes at the deployed cutoff;
\item a nonlinear FK amplitude is not defined without a smoothing scale,
and the formalism does not currently supply one;
\item quoting a nonlinear number therefore requires either adopting a beam
width as part of the observable's definition, or an effective-theory
treatment in which the ultraviolet-sensitive part is absorbed into a
counterterm fixed by simulation or data.
\end{itemize}
